\documentclass[french, 11pt]{article}

\input{setup.tex}

\def\chapitre{10}
\def\pagetitle{Déduction naturelle.}

\begin{document}

\input{title.tex}

Le but de ce chapitre est de commencer à formaliser ce qu'est une preuve. On pourrait ainsi essayer d'en générer, ou d'utiliser ce formalisme pour vérifier des preuves. Une preuve au sens usuel est un raisonnement mathématique qui, à partir d'un ensemble d'hypothèses amène à une conclusion. 

\section{Rappels de logique propositionnelle.}

\begin{rappel}{Définition syntaxique.}{}
    Si $\V$ est un ensemble de variables propositionnelles, on peut définir des formules propositionnelles sur les variables de $\V$ par induction, pour $\phi$ et $\psi$ deux formules propositionnelles :\\
    $\hra$ $\top$ et $\bot$ sont des formules propositionnelles.\\
    $\hra$ $x\in\V$ est une formule propositionnelle.\\
    $\hra$ $\phi \land \psi$ et $\phi \lor \psi$ sont des formules propositionnelles.\\
    $\hra$ $\lnot\phi$ est une formule propositionnelle.\\
    On a défini les formules propositionnelles syntaxiquement : on peut simplement reconnaître les formules propositionnelles.
\end{rappel}

\vspace*{-0.9cm}

\begin{rappel}{Définition sémantique.}{}
    On associe aux formules propositionnelles une définition sémantique : on leur donne un sens.\\
    On définit les valuations : application de $\V$ dans $\{0,1\}$, et on l'étend aux formules propositionnelles inductivement:\\
    $\hra$ Si $\phi=\top$, on pose $v(\phi)=1$ et si $\phi=\bot$, $v(\phi)=0$.\\
    $\hra$ Si $\phi=x\in\V$, on pose $v(\phi)=v(x)$.\\
    $\hra$ Si $\phi=\phi_1\land\phi_2$, on pose $v(\phi)= v(\phi_1)v(\phi_2)$.\\
    $\hra$ Si $\phi=\phi_1\lor\phi_2$, on pose $v(\phi)=v(\phi_1)+v(\phi_2) - v(\phi_1)v(\phi_2)$.\\
    $\hra$ Si $\phi=\lnot\phi_1$, on pose $v(\phi) = 1 - v(\phi_1)$.
\end{rappel}

\vspace*{-0.9cm}

\begin{defi}{Modèle.}{}
    On note $\M(\phi)=\{v:\V\to\{0,1\}, ~ v(\phi)=1\}$, appelé \bf{modèle} de $\phi$. On peut le construire inductivement :\\
    $\hra$ $\M(\bot)=\0$ et $\M(\top)=\{v:\V\to\{0,1\}\}$.\\
    $\hra$ $\M(x)=\{v:\V\to\{0,1\}, ~ v(x)=1\}$.\\
    $\hra$ $\M(\phi_1\land\phi_2)=\M(\phi_1)\cap\M(\phi_2)$.\\
    $\hra$ $\M(\phi_1\lor\phi_2)=\M(\phi_1)\cup\M(\phi_2)$.\\
    $\hra$ $\M(\lnot\phi_1)=\ov{\M(\phi_1)}$.\\
    Si $\M(\phi)\neq\0$, on dit que $\phi$ est \bf{satisfiable}; si $\M(\phi)=\{v:\V\to\{0,1\}\}$, $\phi$ est une \bf{tautologie}.\\
    On dit que $\phi$ et $\psi$ sont \bf{sémantiquement équivalentes} si $\M(\phi)=\M(\psi)$, on note $\phi\equiv\psi$.
\end{defi}

\begin{ex}{}{}
    Puisque $\forall \phi,\psi ~ : ~ \left[(\phi \ra \psi) \equiv (\lnot \phi \lor \psi)\right]$, on peut montrer par induction que pour toute formule $\phi$, il existe une formule $\phi'$ telle que $\phi'\equiv \phi$ et n'utilisant pas le symbole $\ra$.
\end{ex}

\vspace*{-0.8cm}

\begin{ex}{}{}
    Puisque $\lnot(\lnot \phi \land \lnot \psi) \equiv (\phi \lor \psi)$, on peut démontrer que seuls $\land$ et $\lnot$ sont suffisants pour décrire toutes les formules propositionnelles (ou $\lor$ et $\lnot$).
\end{ex}

Lorsqu'on prouve quelque chose, on cherche à prouver que la formule sous-jacente est une tautologie.\\
-- On peut tester chaque valuation (complexité exponentielle.)

\section{Déduction naturelle.}

\begin{defi}{}{}
    On appelle séquent un couple $(\G,\phi)$ où $\G$ est un ensemble de formules et $\phi$ une formule.\\
    On le note $\G\vdash\phi$, et si $\G=\0$, on notera $\vdash\phi$.
\end{defi}

\vspace*{-0.8cm}

\begin{defi}{}{}
    On appelle règle d'inférence un $(n+1)$-uplet de séquents et on la note : 
    \begin{center}
        \begin{prooftree}
            \AxiomC{$\G_1\vdash\phi_1$}
            \AxiomC{$\G_2\vdash\phi_2$}
            \AxiomC{$\cdots$}
            \AxiomC{$\G_n\vdash\phi_n$}
            \QuaternaryInfC{$\G\vdash\phi$}
        \end{prooftree}
    \end{center}
    Les $n$ premiers séquents sont les prémisses, le dernier est la conclusion.
\end{defi}

\vspace*{-0.8cm}

\begin{ex}{Modus Ponens.}{}
    On appelle élimination de l'implication la règle d'inférence :
    \begin{prooftree}
        \AxiomC{$\G\vdash\phi$}
        \AxiomC{$\G\vdash\phi\ra\psi$}
        \BinaryInfC{$\G\vdash\psi$}
    \end{prooftree}
    Si avec les hypothèses $\G$ on prouve $\phi$, et que $\phi$ implique $\psi$, alors on peut prouver $\psi$ sous les hypothèses $\G$.
\end{ex}

\vspace*{-0.8cm}

\begin{ex}{}{}
    On appelle introduction de l'implication :
    \begin{prooftree}
        \AxiomC{$\G, \phi\vdash\psi$}
        \UnaryInfC{$\G\vdash\phi\ra\psi$}
    \end{prooftree}
    Si on arrive à prouver $\psi$ en supposant $\G$ et $\phi$, on arrive à prouver $\phi\ra\psi$ sans supposer $\phi$, seulement $\G$.
\end{ex}

\vspace*{-0.8cm}

\begin{defi}{Arbre de preuve.}{}
    Si $\G\vdash\phi$ est un séquent, on appelle arbre de preuve de $\G\vdash\phi$ :\\
    $\bullet$ Si $\G\vdash\phi$ appartient au systèmé déductif, l'arbre est $\ov{\G\vdash\phi}$.\\
    $\bullet$ Sinon :
    \begin{prooftree}
        \AxiomC{$\G_1\vdash\phi_1$}
        \AxiomC{$\G_2\vdash\phi_2$}
        \AxiomC{$\cdots$}
        \AxiomC{$\G_n\vdash\phi_n$}
        \QuaternaryInfC{$\G\vdash\phi$}
    \end{prooftree}
    C'est une définition par récurrence.
\end{defi}

\begin{ex}{}{}
    Trouvons un arbre de preuve pour $\vdash(A\ra B)\ra(B\ra C)\ra(A\ra C)$.
    \tcblower
    \begin{prooftree}
        \AxiomC{}
        \RightLabel{ax}
        \UnaryInfC{$\G\vdash B\ra C$}
        \AxiomC{}
        \RightLabel{ax}
        \UnaryInfC{$\G\vdash A \ra B$}
        \AxiomC{}
        \RightLabel{ax}
        \UnaryInfC{$\G\vdash A$}
        \RightLabel{$\Rightarrow$e}
        \BinaryInfC{$\G\vdash B$}
        \RightLabel{$\Rightarrow$e}
        \BinaryInfC{$\G:= ~ A\ra B, B\ra C, A\vdash C$}
        \RightLabel{$\Rightarrow$i}
        \UnaryInfC{$A\ra B, B\ra C \vdash A\ra C$}
        \RightLabel{$\Rightarrow$i}
        \UnaryInfC{$A\ra B \vdash (B\ra C)\ra(A\ra C)$}
        \RightLabel{$\Rightarrow$i}
        \UnaryInfC{$\vdash(A\ra B)\ra(B\ra C)\ra(A\ra C)$}
    \end{prooftree}
\end{ex}

\begin{defi}{}{}
    On définit d'autres règles d'inférence pour les autres constructeurs logiques :
    \begin{center}
        \begin{tabular}{| c | c c c |}
            \hline&&&\\
            $\land$ &
            \begin{bprooftree}
                \AxiomC{$\G\vdash\phi$}
                \AxiomC{$\G\vdash\psi$}
                \RightLabel{$\land_{i}$}
                \BinaryInfC{$\G\vdash\phi\land\psi$}
            \end{bprooftree}
            &
            \begin{bprooftree}
                \AxiomC{$\G\vdash\phi\land\phi$}
                \RightLabel{$\land_{eg}$}
                \UnaryInfC{$\G\vdash\phi$}
            \end{bprooftree}
            &
            \begin{bprooftree}
                \AxiomC{$\G\vdash\phi\land\phi$}
                \RightLabel{$\land_{ed}$}
                \UnaryInfC{$\G\vdash\psi$}
            \end{bprooftree}\\&&&\\\hline
        \end{tabular}
    \end{center}
    \vspace*{0.5cm}
    \begin{center}
        \begin{tabular}{|c|c c c c|}
            \hline &&&&\\
            $\lor$ & 
            \begin{bprooftree}
                \AxiomC{$\G\vdash\phi$}
                \RightLabel{$\lor_{ig}$}
                \UnaryInfC{$\G\vdash\phi\lor\psi$}
            \end{bprooftree}
            &
            \begin{bprooftree}
                \AxiomC{$\G\vdash\phi$}
                \RightLabel{$\lor_{id}$}
                \UnaryInfC{$\G\vdash\phi\lor\psi$}
            \end{bprooftree}
            &
            \begin{bprooftree}
                \AxiomC{$\G\vdash\phi\land\phi$}
                \RightLabel{$\land_{ed}$}
                \UnaryInfC{$\G\vdash\psi$}
            \end{bprooftree}
            &
            \begin{bprooftree}
                \AxiomC{$\G\vdash\phi\lor\psi$}
                \AxiomC{$\G,\phi\vdash \a$}
                \AxiomC{$\G,\psi\vdash \a$}
                \RightLabel{$\lor_e$}
                \TrinaryInfC{$\G\vdash\a$}
            \end{bprooftree}\\&&&&\\
            \hline
        \end{tabular}
    \end{center}
        
\end{defi}

\begin{ex}{}{}
    Prouvons $\vdash(A\lor B) \ra (B\lor A)$.
    \begin{prooftree}
        \AxiomC{}
        \RightLabel{ax}
        \UnaryInfC{$A\lor B \vdash A\lor B$}
        \AxiomC{}
        \RightLabel{ax}
        \UnaryInfC{$A\lor B, A\vdash A$}
        \RightLabel{$\lor_{id}$}
        \UnaryInfC{$A\lor B, A\vdash B\lor A$}
        \AxiomC{}
        \RightLabel{ax}
        \UnaryInfC{$A\lor B, B\vdash B$}
        \RightLabel{$\lor_{ig}$}
        \UnaryInfC{$A\lor B, B\vdash B\lor A$}
        \TrinaryInfC{$A\lor B \vdash B \lor A$}
        \RightLabel{$\Rightarrow_i$}
        \UnaryInfC{$\vdash(A\lor B)\ra(B\lor A)$}
    \end{prooftree}
\end{ex}

\begin{ex}{}{}
    Prouvons $\vdash (A\lor B)\land C \ra (A\land C)\lor (B \land C)$.
    \tiny
    \begin{prooftree}
        \AxiomC{}
        \RightLabel{ax}
        \UnaryInfC{$(A\lor B) \land C \vdash (A \lor B) \land C$}
        \UnaryInfC{$(A\lor B)\land C \vdash (A \lor B)$}

        \AxiomC{}
        \RightLabel{ax}
        \UnaryInfC{$(A\lor B)\land C, A \vdash A$}

        \AxiomC{}
        \RightLabel{ax}
        \UnaryInfC{$(A\lor B)\land C, C \vdash (A\lor B)\land C$}
        \RightLabel{$\land_{ed}$}
        \UnaryInfC{$(A\lor B)\land C,A \vdash C$}
        \RightLabel{$\land_{ed}$}
        \BinaryInfC{$(A\lor B)\land C, A \vdash A \land C$}
        \UnaryInfC{$(A\lor B)\land C \vdash C, A\vdash (A\land C)\lor (B\land C)$}
        \AxiomC{$(A\lor B)\land C, B \vdash (A\land C)\lor (B\land C)$}
        \RightLabel{$\land_{ed}$}
        \TrinaryInfC{$(A\lor B)\land C \vdash (A\land C)\lor(B\land C)$}
        \UnaryInfC{$\vdash(A\lor B)\land C \ra (A\land C)\lor(B\land C)$}
    \end{prooftree}
\end{ex}

% $\vdash(A\land B)\lor C \ra (A \lor C)\land (B\lor C)$

\end{document}